\subsection*{MT-M10. A point-jet basis forces a polynomial cluster-product bound}
The counterexample MT-M09 has a quantitative converse obstruction.
Let $\mathcal E$ be the exponential-Schwartz space of MT-M06,
let $J\subset\mathcal E$ be a closed linear subspace, and put
$Q=\mathcal E/J$. Fix a real $\sigma$ and a discrete set
$Z\subset\sigma+i\mathbb R$, with positive integer multiplicities
$m_\rho$. Suppose every functional
\[
 \delta_{\rho,k}(f)=\frac{(\mathcal Mf)^{(k)}(\rho)}{k!},
 \qquad 0\le k<m_\rho,
\]
annihilates $J$ and occurs as a coefficient functional of a
Schauder basis of $Q$ with continuous coordinate functionals.
Other basis coordinates are permitted.
Then there are constants $c>0$ and $A\ge0$, independent of
$\rho=\sigma+it$ and of the finite subset
$W\subset Z\setminus\{\rho\}$, such that
\[
       \prod_{\nu\in W}|t-\operatorname{Im}\nu|^{m_\nu}
                    \ge c(1+|t|)^{-A}.                \tag{MT-M10}
\]
In particular the distinct ordinates have a uniform polynomial
separation bound. The conclusion already follows from an ordinary
Schauder basis; absolute convergence is not needed.

\begin{proof}
The seminorms $p_{N,j}(f)=\sup_x e^{N|x|}|f^{(j)}(x)|$ make
$\mathcal E$ a Fr\'echet space: a Cauchy sequence converges uniformly,
with all derivatives, on every compact interval, and the weighted
uniform limits give the same smooth limit and all the seminorm
limits. The closed quotient $Q$ is Fr\'echet as well.
For a continuous seminorm $p$ on $\mathcal E$, write
\[
 \overline p([f])=\inf_{u\in J}p(f+u).
\]
The quotient seminorms of finite maxima of the $p_{N,j}$ generate
the quotient topology.

First fix any such $\overline p$. Let $S_n$ be the partial-sum
operators for the proposed basis, and let $R_n=S_n-S_{n-1}$ be
its rank-one coordinate projections, with $S_0=0$.
Convergence of every basis expansion makes
$\sup_n\overline p(R_nq)$ finite for each $q\in Q$.
For completeness, the uniform-boundedness step follows directly
from Baire's theorem. The closed sets
\[
 A_k=\{q:\sup_n\overline p(R_nq)\le k\},\qquad k=1,2,\ldots,
\]
cover $Q$. One contains $q_0+U$, where $U$ is a neighborhood of
zero, and $q_0$ itself may be chosen in that set. Subtraction
gives $\overline p(R_nu)\le2k$ for $u\in U$, uniformly in $n$.
Choose $U$ to contain a unit ball for a finite maximum of
quotient seminorms and rescale. This proves that for a continuous
quotient seminorm $\overline q$ and a constant $C$,
\[
             \overline p(R_nq)\le C\overline q(q)
                         \quad\hbox{for every }n,q.    \tag{1}
\]
The same conclusion holds when the seminorm on the right is
zero, by applying the bound to arbitrary scalar multiples.

Let $b_{\rho,0}$ be the basis vector with coefficient
$\delta_{\rho,0}$. For a fixed Gaussian $g_a$ with
$\mathcal Mg_a(s)=e^{as^2}$, use the test function
\[
                f_t(x)=e^{-\sigma x-itx}g_a(x).
\]
It satisfies $\delta_{\rho,0}(f_t)=1$.
Every fixed seminorm of $f_t$ is bounded by a fixed power of
$1+|t|$, by the Leibniz formula and Gaussian decay.
Apply (1) to $[f_t]$ and its coordinate projection. For each
fixed $\overline p$ there are constants $C_p,A_p$ such that
\[
       \overline p(b_{\rho,0})\le C_p(1+|t|)^{A_p}
                                      \quad(\rho\in Z). \tag{2}
\]
The constants do not depend on the enumeration position of
$b_{\rho,0}$.

We next give a uniform estimate for confluent divided differences.
For real nodes $t_0,\ldots,t_m$, repeated nodes allowed, the
divided difference of $u\mapsto e^{ixu}$ is
\[
 [t_0,\ldots,t_m]e^{ixu}
   =\int_{\Delta_m}(ix)^m
           e^{ix\sum_{j=0}^m\lambda_jt_j}\,d\lambda,
 \qquad
 \Delta_m=\{\lambda_j\ge0:\textstyle\sum\lambda_j=1\},
                                                        \tag{3}
\]
where $d\lambda$ has total mass $1/m!$.
One elementary verification suffices also for repeated nodes.
For distinct nodes, the divided difference of $u^k$ is zero
when $k<m$ and is the complete homogeneous polynomial
$h_{k-m}(t_0,\ldots,t_m)$ when $k\ge m$. This follows by
expanding the partial-fraction identity
\[
 \sum_{j=0}^m
 \frac{1}{(1-zt_j)\prod_{\ell\ne j}(t_j-t_\ell)}
          =\frac{z^m}{\prod_{j=0}^m(1-zt_j)}.
\]
On the simplex,
$\int_{\Delta_m}\lambda_0^{a_0}\cdots\lambda_m^{a_m}\,d\lambda
=\prod_j a_j!/(m+\sum a_j)!$; successive integration of
monomials proves this formula. Expanding the exponential and
using this integral gives exactly the same power series,
proving (3). Continuity as nodes coalesce gives the confluent
version, whose functionals are derivatives divided by factorials.
In particular the absolute value in (3) is at most $|x|^m/m!$.

Take $t_0=t$ and repeat each $\operatorname{Im}\nu$, $\nu\in W$,
exactly $m_\nu$ times among the other nodes. Apply the divided
difference to $u\mapsto\mathcal Mf(\sigma+iu)$.
All resulting point-jet functionals annihilate $J$.
Choose an integer $N\ge|\sigma|+1$ and put $p=p_{N,0}$.
Formula (3), Fubini's theorem and exponential decay give
\[
 \left|[t_0,\ldots,t_m]\mathcal Mf(\sigma+iu)\right|
 \le\int_{\mathbb R}\frac{|x|^m}{m!}e^{\sigma x}|f(x)|\,dx
 \le2p(f).                                             \tag{4}
\]
The last constant is independent of both the nodes and $m$,
since $\int_{\mathbb R}|x|^m e^{-|x|}\,dx=2m!$.
The same inequality holds on $Q$ with $\overline p$.

Biorthogonality says that on $b_{\rho,0}$ all the other
point jets in this divided difference vanish. The coefficient
of the simple node $t$ is
$\prod_{\nu\in W}(t-\operatorname{Im}\nu)^{-m_\nu}$.
For repeated other nodes this follows either from partial
fractions or by coalescing distinct nodes away from $t$.
Consequently (4) and (2) imply
\[
 \prod_{\nu\in W}|t-\operatorname{Im}\nu|^{-m_\nu}
       \le2\overline p(b_{\rho,0})
       \le2C_p(1+|t|)^{A_p},
\]
which is (MT-M10).
\end{proof}

For any fixed $0<r<1$, choosing all other zeros within distance
$r$ also bounds their total multiplicity by
$(A\log(1+|t|)-\log c)/|\log r|$.
The product condition is stronger than such a local count:
it controls the joint cost of arbitrarily close clusters.
Taking one other node recovers polynomial separation.
Taking the clusters in MT-M09 contradicts (MT-M10) directly.
This identifies an explicit necessary analytic input for
Masullo 9.3.2.5, PDF~425--426, even after the closed-ideal repair.
It neither asserts that arithmetic zeros violate that condition
nor proves that the condition suffices for a basis.

\subsection*{MT-M11. The same basis would bound arithmetic leading coefficients}
Fix a primitive unitary Hecke $L$-function $L(s,\eta)$ and assume
its nontrivial zeros lie on $\operatorname{Re}s=1/2$.
Suppose a closed quotient as in MT-M10 has all its nontrivial
zero jets as coefficient functionals, with their actual
multiplicities. Then there are constants $c,C,T_0>0$, depending
on the fixed function and the proposed basis, such that at every
nontrivial zero $\rho=1/2+iT$ with $|T|\ge T_0$ and order $m$,
\[
               \left|\frac{L^{(m)}(\rho,\eta)}{m!}\right|
                        \ge c(1+|T|)^{-C}.             \tag{MT-M11}
\]
In particular, at simple zeros the reciprocal of $|L'(\rho,\eta)|$
would have at most polynomial growth. The zero-location input is
$A_D$ for Dirichlet functions and $M_H$ for general Hecke functions;
the asserted quotient basis is an additional hypothesis.

\begin{proof}
We use the fixed-function analytic facts proved in NT-Z10:
polynomial growth on a fixed vertical strip away from its
finitely located poles, and $O(\log(2+|T|))$ zeros, with
multiplicity, in each bounded-height window.
The absolutely convergent Euler product gives
$|L(2+iT,\eta)|\ge c_\eta>0$: each local inverse factor
has modulus at least $(1+(\operatorname N\mathfrak p)^{-2})^{-1}$,
and the product of these lower bounds is positive.

Put $c_T=2+iT$. Choose $4\le R\le5$ so that the boundary of
$D=\{s:|s-c_T|<R\}$ contains no zero. At sufficiently large
$|T|$, it contains no pole or trivial zero either, because
the imaginary parts of their finitely many arithmetic
progressions parallel to the real axis are fixed by the infinity type.
Its zeros are therefore precisely the relevant critical-line
zeros in that disk, say $\rho_j$ with multiplicities $m_j$.
There are
$M=\sum_jm_j=O_\eta(\log(2+|T|))$ of them.
In the normalized coordinate $z=(s-c_T)/R$, put
$a_j=(\rho_j-c_T)/R$ and form the finite Blaschke product
\[
 B(s)=\prod_j
       \left(\frac{z-a_j}{1-\overline{a_j}z}\right)^{m_j}.
\]
Constant unimodular normalizations of the factors are
irrelevant. The function $G=L/B$ is holomorphic and zero-free
on the disk and its boundary. On the boundary $|B|=1$.
The polynomial growth bound and the maximum principle give
$|G(s)|\le M_T=C_1(1+|T|)^{C_2}$ throughout $D$, enlarging
$M_T$ if necessary to make it at least one.
At the center, $|B(c_T)|\le1$, so
$|G(c_T)|\ge c_\eta$.

The nonnegative harmonic function
$u(s)=\log M_T-\log|G(s)|$ satisfies the disk Harnack bound
\[
 u(\rho)\le
 \frac{R+|\rho-c_T|}{R-|\rho-c_T|}\,u(c_T)
 \le\frac{11}{5}\,u(c_T),
\]
because $|\rho-c_T|=3/2$ and $R\ge4$.
This elementary bound follows by comparing the maximum and
minimum of the Poisson kernel on each concentric boundary
circle and then letting that circle increase to radius $R$.
It implies $|G(\rho)|\ge c_2(1+|T|)^{-C_3}$.

Let $a=(\rho-c_T)/R$, so $|a|\le3/8$.
The leading coefficient of its own factor in $B$ contributes
$[R(1-|a|^2)]^{-m}\ge5^{-m}$.
Every other factor at $\rho$ has modulus
\[
 \frac{|a-a_j|}{|1-\overline{a_j}a|}
 \ge\frac{|\rho-\rho_j|}{10}.
\]
Thus, with $W$ the distinct other zeros in $D$,
\[
 \left|\lim_{s\to\rho}\frac{B(s)}{(s-\rho)^m}\right|
 \ge 10^{-M}\prod_{\rho_j\in W}|\rho-\rho_j|^{m_j}.
\]
All these points are on the same vertical line. Apply
(MT-M10) to this finite set $W$. Since
$10^{-M}\ge c_3(1+|T|)^{-C_4}$, the displayed expression
has a polynomial lower bound. Multiplying by the lower bound
for $|G(\rho)|$ proves (MT-M11).
\end{proof}

This is a necessary consequence of the proposed arithmetic
point-jet basis, not a derivative estimate obtained from GRH
alone. It converts the remaining interpolation problem into
a concrete additional test on arithmetic leading coefficients.
The work placement is the Masullo 9.3.2.5 repair, NT-Z10,
and the cote~1 trace/Fredholm programme reported in SA23
geo-H-lidskii/geo-H-genus. Positivity of the trace cannot
replace this topological prerequisite.
